\section{Introduzione}
L’obiettivo di questo progetto è creare una fft a 16 campioni applicando l’algoritmo di Cooley-Tukey, che permette di semplificare e ridurre il numero delle operazioni necessarie. Date quindi le richieste e le limitazioni sui componenti da impiegare dovevamo ottenere le migliori performance possibili.
Si sono quindi sviluppate 2 modalità d’uso (descritte nella sottosezione \ref{subsec:modalità}): singola e continua. La seconda permette di aumentare il throughput a una solve ogni 6 colpi di clock.
